\section{BÀI 8: PHÂN TÍCH THỂ TÍCH}
\subsection{THÍ NGHIỆM}
\begin{enumerate}
    \item \textbf{Thí nghiệm 1}
    \begin{itemize}
        \item Xây dựng đường cong chuẩn độ một axit mạnh bằng một bazơ mạnh dựa theo bảng.
    \end{itemize}
    \item \textbf{Thí nghiệm 2}
    \begin{itemize}
        \item Chuẩn độ axit – bazo với thuốc thử phenolphtalein
        \item Tráng buret bằng dung dịch NaOH 0.1N, cho từ từ dung dịch NaOH 0.1N vào buret
        \item Dùng pipet 10ml lấy 10ml dd HCl cho vào erlen 150ml, thêm 10ml nước cất và 2 giọt phenolphtalein
        \item Bắt đầu tiến hành chuẩn độ
        \item Lặp lại thí nghiệm 3 lần để tính giá trị trung bình.
    \end{itemize}
    \item \textbf{Thí nghiệm 3}
    \begin{itemize}
        \item Tương tự thí nghiệm 2 nhưng thay phenolphtalein bằng metyl da cam. Màu dung dịch đổi từ đỏ sang cam.
    \end{itemize}
    \item \textbf{Thí nghiệm 4}
    \begin{itemize}
        \item Tương tự thí nghiệm 2 nhưng thay dung dịch HCl bằng dung dịch axit axetic.
        \item Bởi vì axit acetic là một axit yếu nên không thể dùng chất chỉ thị metyl da cam.
        \item Làm thí nghiệm 3 lần nhưng chỉ với chất chỉ thị phenolphtalein.
    \end{itemize}
\end{enumerate}

\subsection{Kết quả thí nghiệm}
\subsubsection{Thí nghiệm 1}
\begin{itemize}
    \item Xác định:
    \begin{itemize}
        \item pH điểm tương đương: 7.26.
        \item Bước nhảy pH: Từ pH 3.36 đến pH 10.56.
        \item Chất chỉ thị thích hợp: Phenolphthalein, metyl da cam.
    \end{itemize}
\end{itemize}

\subsubsection{Thí nghiệm 2}
\begin{itemize}
    \item Tráng buret bằng dung dịch NaOH 0.1N, sau đó cho từ từ dung dịch NaOH 0.1N vào buret. Chỉnh mức dung dịch ngang vạch 0.
    \item Dùng pipet 10 ml lấy 10 ml dung dịch HCl chưa biết nồng độ vào erlen 150ml, thêm 10 ml nước cất và 2 giọt phenolphtalein.
    \item Mở khóa buret nhỏ từ từ dung dịch NaOH xuống erlen, vừa nhỏ vừa lắc nhẹ đến khi dung dịch trong erlen chuyển sang màu hồng nhạt bền thì khóa buret. Đọc thể tích dung dịch NaOH đã dùng.
    \item Lặp lại thí nghiệm hai lần nữa để tính giá trị trung bình.
    \item Điền đầy đủ các giá trị vào bảng sau:
    \begin{table}[H]
        \centering
        \begin{tabular}{|c|c|c|c|c|c|}
            \hline
            \textbf{Lần} & \textbf{V\textsubscript{HCl} (ml)} & \textbf{V\textsubscript{NaOH} (ml)} & \textbf{C\textsubscript{NaOH} (N)} & \textbf{C\textsubscript{HCl} (N)} & \textbf{Sai số} \\
            \hline
            1 & 10 & 10 & 0.1 & 0.1 & 0.005 \\
            \hline
            2 & 10 & 10.1 & 0.1 & 0.101 & 0.005 \\
            \hline
        \end{tabular}
    \end{table}
    \item $\overline{C_{HCl}} = \frac{1}{2} \sum_{i=1}^{2} C_{HCl\_i} = \frac{1}{2}(0.1 + 0.101) = 0.1005 \text{ (N)}$
    \item Ta tính được sai số tuyệt đối của mỗi lần đo theo công thức: $\Delta C_{HCl} = |C_{HCl\_i} - \overline{C_{HCl}}|$
    \item => Các giá trị sai số được thể hiện ở bảng trên.
    \item Sai số trung bình: $\overline{\Delta C_{HCl}} = \frac{1}{2} \sum_{i=1}^{2} \Delta C_{HCl\_i} = \frac{1}{2}(0.005 + 0.005) = 0.005 \text{ (N)}$
    \item Suy ra, nồng độ HCl cần dùng là: $C_{HCl} = 0.1005 \pm 0.005 \text{ (N)}$
\end{itemize}

\subsubsection{Thí nghiệm 3}
\begin{itemize}
    \item Tiến hành tương tự thí nghiệm 2 nhưng thay chất chỉ thị phenolphtalein bằng metyl da cam. Màu dung dịch đổi từ đỏ sang cam.
    \item Điền đầy đủ các giá trị vào bảng sau:
    \begin{table}[H]
        \centering
        \begin{tabular}{|c|c|c|c|c|c|}
            \hline
            \textbf{Lần} & \textbf{V\textsubscript{HCl} (ml)} & \textbf{V\textsubscript{NaOH} (ml)} & \textbf{C\textsubscript{NaOH} (N)} & \textbf{C\textsubscript{HCl} (N)} & \textbf{Sai số} \\
            \hline
            1 & 10 & 10.2 & 0.1 & 0.102 & 0.0015 \\
            \hline
            2 & 10 & 9.9 & 0.1 & 0.099 & 0.0015 \\
            \hline
        \end{tabular}
    \end{table}
    \item $\overline{C_{HCl}} = \frac{1}{2} \sum_{i=1}^{2} C_{HCl\_i} = \frac{1}{2}(0.102 + 0.099) = 0.1005 \text{ (N)}$
    \item Ta tính được sai số tuyệt đối của mỗi lần đo theo công thức: $\Delta C_{HCl} = |C_{HCl\_i} - \overline{C_{HCl}}|$
    \item => Các giá trị sai số được thể hiện ở bảng trên.
    \item Sai số trung bình: $\overline{\Delta C_{HCl}} = \frac{1}{2} \sum_{i=1}^{2} \Delta C_{HCl\_i} = \frac{1}{2}(0.0015 + 0.0015) = 0.0015 \text{ (N)}$
    \item Suy ra, nồng độ HCl cần dùng là: $C_{HCl} = 0.1005 \pm 0.0015 \text{ (N)}$
\end{itemize}

\subsubsection{Thí nghiệm 4}
\begin{itemize}
    \item Tiến hành tương tự thí nghiệm 2 nhưng thay chất dung dịch HCl bằng dung dịch axit axetic. Làm thí nghiệm 2 lần với lần đầu dùng chất chỉ thị phenolphtalein, lần sau dùng metyl da cam.
\end{itemize}
\textbf{a. Thí nghiệm 4a}
\begin{itemize}
    \item Điền đầy đủ các giá trị vào bảng sau:
    \begin{table}[H]
        \centering
        \begin{tabular}{|c|c|c|c|c|c|}
            \hline
            \textbf{Lần} & \textbf{V\textsubscript{CH\textsubscript{3}COOH} (ml)} & \textbf{V\textsubscript{NaOH} (ml)} & \textbf{C\textsubscript{NaOH} (N)} & \textbf{C\textsubscript{HCl} (N)} & \textbf{Sai số} \\
            \hline
            1 & 10 & 9 & 0.1 & 0.09 & 0.0005 \\
            \hline
            2 & 10 & 9.1 & 0.1 & 0.091 & 0.0005 \\
            \hline
        \end{tabular}
    \end{table}
    \item $\overline{C_{CH_3COOH}} = \frac{1}{2} \sum_{i=1}^{2} C_{CH_3COOH\_i} = \frac{1}{2}(0.09 + 0.091) = 0.0905 \text{ (N)}$
    \item Ta tính được sai số tuyệt đối của mỗi lần đo theo công thức: $\Delta C_{CH_3COOH} = |C_{CH_3COOH\_i} - \overline{C_{CH_3COOH}}|$
    \item => Các giá trị sai số được thể hiện ở bảng trên.
    \item Sai số trung bình: $\overline{\Delta C_{CH_3COOH}} = \frac{1}{2} \sum_{i=1}^{2} \Delta C_{CH_3COOH\_i} = \frac{1}{2}(0.0005 + 0.0005) \approx 0.0005 \text{ (N)}$
    \item Suy ra, nồng độ HCl cần dùng là: $C_{CH_3COOH} = 0.0905 \pm 0.0005 \text{ (N)}$
\end{itemize}

\textbf{b. Thí nghiệm 4b}
\begin{itemize}
    \item Tiến hành tương tự thí nghiệm 2 nhưng thay chất chỉ thị phenolphtalein bằng metyl da cam. Màu dung dịch đổi từ đỏ sang cam.
    \item Điền đầy đủ các giá trị vào bảng sau:
    \begin{table}[H]
        \centering
        \begin{tabular}{|c|c|c|c|c|c|}
            \hline
            \textbf{Lần} & \textbf{V\textsubscript{CH\textsubscript{3}COOH} (ml)} & \textbf{V\textsubscript{NaOH} (ml)} & \textbf{C\textsubscript{NaOH} (N)} & \textbf{C\textsubscript{HCl} (N)} & \textbf{Sai số} \\
            \hline
            1 & 10 & 0.9 & 0.1 & 0.009 & 0.0005 \\
            \hline
            2 & 10 & 1 & 0.1 & 0.01 & 0.0005 \\
            \hline
        \end{tabular}
    \end{table}
    \item $\overline{C_{CH_3COOH}} = \frac{1}{2} \sum_{i=1}^{2} C_{CH_3COOH\_i} = \frac{1}{2}(0.009 + 0.01) = 0.0095 \text{ (N)}$
    \item => Vì axit acetic là một axit yếu nên khi sử dụng số liệu thí nghiệm với chất chỉ thị metyl da cam sẽ cho kết quả sai, vì vậy chúng ta bỏ qua trường hợp này.
\end{itemize}

\subsection{TRẢ LỜI CÂU HỎI}
\begin{enumerate}
    \item \textbf{Khi thay đổi nồng độ HCl và NaOH, đường cong chuẩn độ có thay đổi hay không? Tại sao?}
    \begin{itemize}
        \item Khi thay đổi nồng độ HCl và NaOH, đường cong chuẩn độ không thay đổi do phương pháp chuẩn độ HCl bằng NaOH được xác định dựa trên phương trình:
        \begin{align*}
            &\text{HCl} + \text{NaOH} \rightarrow \text{NaCl} + \text{H}_2\text{O} \\
            &C_{HCl} \cdot V_{HCl} = C_{NaOH} \cdot V_{NaOH}
        \end{align*}
        \item Với $V_{HCl}$ và $C_{NaOH}$ cố định nên khi $C_{HCl}$ tăng hay giảm thì $V_{NaOH}$ cũng tăng hay giảm theo. Cho nên dù mở rộng ra hay thu hẹp lại thì đường cong chuẩn độ không đổi.
        \item Tương tự đối với trường hợp thay đổi nồng độ NaOH.
    \end{itemize}
    \item \textbf{Việc xác định nồng độ axit HCl trong các thí nghiệm 2 và 3 cho kết quả nào chính xác hơn? Tại sao?}
    Xác định nồng độ axit HCl trong thí nghiệm 2 cho kết quả chính xác hơn. Vì phenol phtalein giúp ta xác định màu chính xác hơn, rõ ràng hơn, do chuyển từ không màu sang hồng nhạt và dễ nhận thấy hơn từ màu đỏ sang da cam.
    \item \textbf{Từ kết quả thí nghiệm 4, việc xác định nồng độ dung dịch axit axetic bằng chất chỉ thị màu nào chính xác hơn? Tại sao?}
    Phenolphtalein chính xác hơn metyl da cam vì axit axetic là axit yếu với điểm định mức lớn hơn 7 nên dùng phenolphtalein sẽ chính xác hơn metyl da cam (bước nhảy 3.0 – 4.4 quá lớn nên không chính xác).
    \item \textbf{Trong phép phân tích thể tích, nếu đổi vị trí của NaOH và axit thì kết quả có thay đổi không? Tại sao?}
    Không thay đổi vì trường hợp này cũng chỉ là phản ứng cân bằng.
\end{enumerate}
